\par\medskip
\begin{center}
\textls[100]{e) Die Tagebuchnotizen des Jahres 1812. \textsc{Gauss}sche Logarithmen.}
\end{center}
\noindent\hspace*{\parindent}In das seit der Notiz 140. vom 29. Juni 1809 nicht benutzte \textit{Tagebuch}
schreibt \textsc{Gauss} am 29. Februar 1812: »Das vorstehende, in Folge von Schicksalsschlägen
abermals unterbrochene Verzeichnis nehmen wir zu Beginn des
Jahres 1812 wieder auf. Im Monat November 1811 gelang es, den rein analytischen
Beweis des Fundamentalsatzes der Lehre von den Gleichungen zu
vervollständigen; da aber nichts zu Papier gebracht worden war, entschwand
ein wesentlicher Teil völlig dem Gedächtnis. Diesen, während einer recht langen
Zeit vergebens gesuchten Teil, haben wir nunmehr glücklich wiedergefunden.«
\par
Gemeint ist hier der zweite, am 7. Dezember 1815 der Gesellschaft der
Wissenschaften vorgelegte Beweis für die Wurzelexistenz der algebraischen
Gleichungen\footnote{\textit{Demonstratio nova altera etc.}, Comment. soc. reg. scient. Gott. rec. vol. III, 1816, Werke III, S. 31, \textit{Anzeige} ebenda, S. 105.}, in bezug auf den \textsc{Gauss} am 19. Februar 1826 an \textsc{Olbers}
schreibt\footnote{\textit{Briefwechsel Olbers-Gauss} II, S. 439.}: »Ich habe in meinem wissenschaftlichen Leben öfter den Fall
gehabt, dass ich durch äussere Umstände veranlasst, Beschäftigungen, die nicht
glückten, bei Seite legte, und die allerdings später glückten, z. B. mein Beweis
für das Haupttheorem der Lehre von den Gleichungen, der in dem 3.
Bande unserer Comm. steht; aber ich habe nachher die zehnfache Anstrengung
gehabt, nur erst wieder auf den Punkt zu kommen, auf dem ich schon
früher mehr als einmal gewesen war.«
\par
Die beiden folgenden Tagebucheintragungen vom 26. September und
15. Oktober 1812 beziehen sich auf die Attraktion der elliptischen Sphäroide
und berichten von der Entdeckung einer durchaus neuen Theorie dieser Attraktion
und von einer »alles übrige erledigenden neuen Methode von wunderbarer
Einfachheit«. Die betreffende Abhandlung wurde der G. d. W. am
28. März 1813 vorgelegt\footnote{\textit{Theoria attractionis corporum sphaeroidicorum ellipticorum homogeneorum}, Comm. soc. reg. scient. Gotting. rec. vol. II, 1813, Werke V, S. 1, \textit{Anzeige} ebenda, S. 279. Vergl. \textsc{Galle}, Werke XI, 2, Abhandlung 1, S. 43.}, aber schon am 5. November 1812 legt \textsc{Gauss}\par
\clearpage\SourceLeaf{240}\ArithmeticHead{155}{R}{VIIIa. DIE ABHANDLUNGEN VON 1813 bis 1816.}
\noindent einem Briefe an \textsc{Laplace}\footnote{Werke X, 1, S. 378, 379, \textsc{Laplaces} Antwort ebenda, S. 380, 381.} einen »Extrait« derselben bei\footnote{Abgedruckt Werke XII, S. 110—113; vergl. die Bemerkungen ebenda, S. 114, wo über die wechselvollen Schicksale dieser Handschrift berichtet wird, die sich zur Zeit in der Bibliothek der \textsc{Mittag-Leffler}-Stiftung zu Djursholm befindet.}, mit der Bitte,
ihn dem Institut de France mitzuteilen; »Sie werden mit Vergnügen sehen« —
schreibt \textsc{Gauss} — »dass mir zwei Seiten genügen, um die vollständige Lösung zu
geben«. Die Besprechung der beiden genannten Abhandlungen bringt der folgende
Abschnitt in der Reihenfolge ihrer Veröffentlichung, hier haben wir nur
noch zu erwähnen, dass \textsc{Gauss} im November 1812 seine \textit{Tafeln zur bequemen Berechnung
der Logarithmen der Summe oder Differenz zweier Grössen, welche selbst
durch ihre Logarithmen gegeben sind} in \textsc{Zachs} Monatlicher Korrespondenz (26,
S. 498) veröffentlicht hat\footnote{Die Einleitung ist abgedruckt Werke III, S. 244—246.}. Wie \textsc{Gauss} selbst hervorhebt, rührt die Idee,
die diesen Tafeln zugrunde liegt, von Z. \textsc{Leonelli} her, dessen 1802—3 in
französischer Sprache\footnote{Z. \textsc{Leonelli}, \textit{Supplément logarithmique etc.}, Bordeaux, an XI.}, 1806 in deutscher Übersetzung erschienene Schrift
\textsc{Gauss} 1808 ausführlich und kritisch besprochen hat\footnote{Allgemeine Literaturzeitung, 1808, S. 353, Werke VIII, S. 121—127; näheres über diese sogenannten \textsc{Gauss}schen Logarithmen findet man bei R. \textsc{Mehmke}, Encyklopädie, Band I, 2, S. 998ff.}.\par
